\documentclass[11pt,a4paper]{article}

\usepackage[T1]{fontenc}
\usepackage[utf8]{inputenc}
\usepackage{lmodern}
\usepackage{amssymb,amsmath,amsthm,mathtools}
\usepackage[margin=2.5cm]{geometry}
\usepackage{parskip}
\usepackage[hidelinks]{hyperref}
\usepackage{microtype}

% Theorem environments (arXiv / journal style)
\theoremstyle{plain}
\newtheorem{theorem}{Theorem}[section]
\newtheorem{proposition}[theorem]{Proposition}
\newtheorem{lemma}[theorem]{Lemma}
\newtheorem{corollary}[theorem]{Corollary}

\theoremstyle{definition}
\newtheorem{definition}[theorem]{Definition}
\newtheorem{example}[theorem]{Example}
\newtheorem{remark}[theorem]{Remark}

% Notation
\newcommand{\Type}{\mathbf{Type}}
\newcommand{\id}{\mathrm{id}}

\title{Associativity of the Zappa--Sz\'{e}p Product\\
  and the Four Cocycle Axioms}
\author{Extracted from a machine-checked development}
\date{}

\begin{document}
\maketitle

\begin{abstract}
We formalise the algebraic equivalence between associativity of the
Zappa--Sz\'{e}p product of two monoids and the four Brin/Zappa--Sz\'{e}p
cocycle axioms. Given monoids $L$ and $G$ with a left action
$\triangleright : G \to L \to L$ and a restriction
$\triangleleft : G \to L \to G$, subject to unit-compatibility conditions that
make $(1_L, 1_G)$ a two-sided unit, the product
\[
  (a,g) * (b,h)
  \;=\;
  \bigl( a \cdot (g \triangleright b),\;
         (g \triangleleft b) \cdot h \bigr)
\]
on $L \times G$ is associative if and only if the four cocycle axioms ZS1--ZS4
hold. The proofs of all stated results have been machine-checked; we record
only the mathematical content.
\end{abstract}

\tableofcontents
\bigskip

%======================================================================
\section{Introduction}
%======================================================================

The Zappa--Sz\'{e}p product is a generalisation of the semidirect product in
which both factors act on each other. Algebraically, one is given two monoids
$L$ and $G$, a left action of $G$ on $L$, and a restriction of $G$ along
elements of~$L$; the product monoid structure on the cartesian product
$L \times G$ is then defined by a mixed formula that uses both operations. For
this product to be associative, the actions must satisfy four cocycle
identities, often called the Brin or Zappa--Sz\'{e}p axioms.

This note records a pure equational proof that, once unit-compatibility is
assumed (so that $(1_L, 1_G)$ is a two-sided unit), associativity of the
product is equivalent to the four cocycle axioms. The forward direction uses
no unit laws. The converse recovers each axiom by specialising associativity
at units and reading off one coordinate: the $L$-coordinate yields ZS1 and
ZS2; the $G$-coordinate yields ZS3 and ZS4.

We work with types as carriers and ordinary functions as operations. Proofs
are indicated by the phrase ``by lean.''

\paragraph{What this note claims.}
For a matched pair of monoids (two monoids, two mutual actions, and four unit
compatibilities),
\[
  \text{ZS1} \wedge \text{ZS2} \wedge \text{ZS3} \wedge \text{ZS4}
  \;\Longleftrightarrow\;
  \text{associativity of the Zappa--Sz\'{e}p product}.
\]
There is nothing more to the associativity law than these four cocycles, and
conversely each cocycle is an instance of associativity.

%======================================================================
\section{Matched pairs of monoids}
%======================================================================

\begin{definition}[Matched pair]\label{def:matched-pair}
A \textbf{matched pair} of monoids consists of:
\begin{enumerate}
  \item monoids $(L, \cdot_L, 1_L)$ and $(G, \cdot_G, 1_G)$, each associative
        and unital;
  \item a left action $\triangleright : G \to L \to L$;
  \item a restriction $\triangleleft : G \to L \to G$;
  \item the four \textbf{unit-compatibility} conditions
        \begin{align}
          \tag{U1a}
          1_G \triangleright a &\;=\; a, \\
          \tag{U1b}
          1_G \triangleleft a &\;=\; 1_G, \\
          \tag{U2a}
          g \triangleright 1_L &\;=\; 1_L, \\
          \tag{U2b}
          g \triangleleft 1_L &\;=\; g,
        \end{align}
        for all $a \in L$ and $g \in G$.
\end{enumerate}
No cocycle axiom is assumed as part of this definition; ZS1--ZS4 are stated
separately below.
\end{definition}

\begin{remark}
The unit-compatibility conditions are exactly what make $(1_L, 1_G)$ a
two-sided unit for the product operation of the next section.
\end{remark}

\begin{definition}[Zappa--Sz\'{e}p product]\label{def:zs-product}
Let $P$ be a matched pair. The \textbf{Zappa--Sz\'{e}p product} is the binary
operation on $L \times G$ defined by
\begin{equation}
  (a,g) * (b,h)
  \;=\;
  \bigl(
    a \cdot_L (g \triangleright b),\;
    (g \triangleleft b) \cdot_G h
  \bigr).
\end{equation}
\end{definition}

%======================================================================
\section{The four cocycle axioms}
%======================================================================

\begin{definition}[Cocycle axioms]\label{def:cocycles}
Let $P$ be a matched pair. The four \textbf{Brin/Zappa--Sz\'{e}p cocycle
axioms} are the following predicates on~$P$.
\begin{align}
  \tag{ZS1}
  g \triangleright (a \cdot_L b)
  &\;=\;
  (g \triangleright a)
    \cdot_L
  \bigl((g \triangleleft a) \triangleright b\bigr),
  \\[6pt]
  \tag{ZS2}
  (g \cdot_G h) \triangleright a
  &\;=\;
  g \triangleright (h \triangleright a),
  \\[6pt]
  \tag{ZS3}
  (g \cdot_G h) \triangleleft a
  &\;=\;
  \bigl(g \triangleleft (h \triangleright a)\bigr)
    \cdot_G
  (h \triangleleft a),
  \\[6pt]
  \tag{ZS4}
  g \triangleleft (a \cdot_L b)
  &\;=\;
  (g \triangleleft a) \triangleleft b.
\end{align}
In words: ZS1 says that $\triangleright$ is a crossed homomorphism in its
$L$-argument; ZS2 says that $\triangleright$ is a left action of~$G$; ZS3
says that $\triangleleft$ is a crossed homomorphism in its $G$-argument; ZS4
says that $\triangleleft$ is a right action of~$L$.
\end{definition}

\begin{definition}[Associativity]\label{def:assoc}
The product of Definition~\ref{def:zs-product} is \textbf{associative} when
\begin{equation}
  \bigl((x * y) * z\bigr) \;=\; \bigl(x * (y * z)\bigr)
\end{equation}
for all $x,y,z \in L \times G$.
\end{definition}

%======================================================================
\section{Cocycles imply associativity}
%======================================================================

\begin{theorem}[Forward direction]\label{thm:assoc-of-zs}
Let $P$ be a matched pair. If ZS1--ZS4 all hold, then the Zappa--Sz\'{e}p
product is associative.
\end{theorem}

\begin{proof}
By lean.
\end{proof}

\begin{remark}[Idea of the argument]
Expand both sides of the associativity equation at a general triple
$((a,g),(b,h),(c,k))$. The $L$-coordinate of the resulting equality follows
from reassociating $\cdot_L$, then ZS1 (to distribute the action over a
product) and ZS2 (the action law). The $G$-coordinate follows from ZS3, then
reassociating $\cdot_G$, then ZS4. No unit law is used.
\end{remark}

%======================================================================
\section{Associativity implies the cocycles}
%======================================================================

\begin{theorem}[Converse direction]\label{thm:zs-of-assoc}
Let $P$ be a matched pair. If the Zappa--Sz\'{e}p product is associative, then
ZS1, ZS2, ZS3, and ZS4 all hold.
\end{theorem}

\begin{proof}
By lean.
\end{proof}

\begin{remark}[How each axiom is recovered]
Each cocycle is an instance of associativity at units.
\begin{enumerate}
  \item \textbf{ZS1.} Specialise at
        $x = (1_L, g)$, $y = (a, 1_G)$, $z = (b, 1_G)$ and read the
        $L$-coordinate; unit laws collapse the remaining factors.
  \item \textbf{ZS2.} Specialise at
        $x = (1_L, g)$, $y = (1_L, h)$, $z = (a, 1_G)$ and read the
        $L$-coordinate.
  \item \textbf{ZS3.} The same triple as for ZS2, reading the $G$-coordinate.
  \item \textbf{ZS4.} The same triple as for ZS1, reading the $G$-coordinate.
\end{enumerate}
The unit-compatibility fields of the matched pair are exactly what make these
specialisations yield the bare cocycle equations.
\end{remark}

%======================================================================
\section{The equivalence}
%======================================================================

\begin{theorem}[Main theorem]\label{thm:zs-iff-assoc}
For any matched pair of monoids, the Zappa--Sz\'{e}p product is associative if
and only if the four cocycle axioms ZS1--ZS4 hold:
\begin{equation}
  \bigl(\mathrm{ZS1} \wedge \mathrm{ZS2}
        \wedge \mathrm{ZS3} \wedge \mathrm{ZS4}\bigr)
  \;\Longleftrightarrow\;
  \mathrm{Assoc}.
\end{equation}
\end{theorem}

\begin{proof}
By lean.
\end{proof}

%======================================================================
\section{An example: the direct product}
%======================================================================

\begin{example}[Trivial actions on $(\mathbb{N}, +, 0)$]\label{ex:nat-direct}
Take $L = G = \mathbb{N}$ with addition as multiplication and $0$ as unit, and
define the trivial actions
\begin{equation}
  g \triangleright a \;=\; a,
  \qquad
  g \triangleleft a \;=\; g.
\end{equation}
These satisfy the four unit-compatibility conditions, so one obtains a matched
pair. All four cocycle axioms hold (both sides of each equation reduce to the
same expression by the definitions of the actions). Consequently, by
Theorem~\ref{thm:zs-iff-assoc}, the Zappa--Sz\'{e}p product is associative. In
this case the product is simply the direct product monoid
$(\mathbb{N} \times \mathbb{N},\, +,\, (0,0))$.
\end{example}

\begin{remark}
The example serves as a sanity check that the matched-pair structure is
inhabited and that ZS1--ZS4 are satisfiable.
\end{remark}

%======================================================================
\section{Discussion}
%======================================================================

\paragraph{What has been established.}
Associativity of the Zappa--Sz\'{e}p product on a matched pair is equivalent to
four equational cocycle axioms. The proof is purely equational: the forward
direction is a direct expansion; the converse is specialisation at units.
Nothing beyond monoid laws, unit compatibility, and the two mutual actions is
used.

\paragraph{Scope.}
The development treats monoids (not necessarily groups). It does not construct
the free or universal matched pair, nor does it address topological, ordered,
or higher-categorical refinements of the Zappa--Sz\'{e}p product. The unit
laws of the product monoid are not stated as separate theorems; they are the
content of the unit-compatibility package built into
Definition~\ref{def:matched-pair}.

\paragraph{Machine checking.}
Every definition and theorem statement above corresponds to a declaration in a
machine-checked formalisation. The phrase ``by lean'' indicates that the proof
has been verified by a proof assistant.

%======================================================================
\section*{Acknowledgements}
%======================================================================

The mathematical content follows the classical theory of Zappa--Sz\'{e}p
products of monoids and the Brin cocycle presentation of the associativity
criteria. The presentation here is a prose rendering of a formal development
intended for human review of the theorem statements.

\begin{thebibliography}{9}

\bibitem{brin2005}
M.~G.~Brin.
\newblock On the Zappa--Sz\'{e}p product.
\newblock \emph{Communications in Algebra}, 33(1):293--304, 2005.

\bibitem{zappa1942}
G.~Zappa.
\newblock Sulla costruzione dei gruppi prodotto di due dati sottogruppi
  permutabili tra loro.
\newblock In \emph{Atti del Secondo Congresso dell'Unione Matematica
  Italiana}, pages 119--125, 1942.

\end{thebibliography}

\end{document}
